% !PRINTED 722
\weberpage{722}

%<<CONT
Reihe der positiven ungeraden Zahlen von 1 bis $\nu - 1$ durchl\"auft;
dann durchl\"auft das ungerade $s$ die vier Zahlenreihen:
\[
t, \quad t+\mu, \quad m-t, \quad m-t-\mu,
\]
und es ist:
\[
\begin{aligned}
\chi(t+\mu) &= -\chi(t), \\
\chi_1(m-t) &= \chi_1(t), \qquad \chi_1(m-\mu-t) = -\chi_1(t), \\
\chi_2(m-t) &= -\chi_2(t), \qquad \chi_2(m-\mu-t) = \chi_2(t).
\end{aligned}
\]
Nach (12) ergiebt sich hieraus:
\[
\begin{aligned}
f(r^{t+\mu}) &= -f(r^t), \\
f_1(r^{m-t}) &= f_1(r^t), \qquad f_1(r^{m-\mu-t}) = -f_1(r^t), \\
f_2(r^{m-t}) &= -f_2(r^t), \qquad f_2(r^{m-\mu-t}) = f_2(r^t).
\end{aligned}
\]
Theilt man demnach die Summen (9) in je vier Theile und benutzt noch
die trigonometrische Gleichung
\[
\sin \frac{(t+\mu)\pi}{m} = \cos \frac{t\pi}{m},
\]
so ergiebt sich
\[
\begin{aligned}
X_1 &= \frac{-1}{m} \sum^{t} f_1(r^t)
\log\left(\operatorname{tg} \frac{t\pi}{m}\right)^2, \\
X_2 &= \frac{i\pi}{m} \sum^{t} f_2(r^t).
\end{aligned}
\]
In dieser Summe ist
\[
\frac{t\pi}{m} < \frac{\pi}{4}, \qquad
\operatorname{tg} \frac{t\pi}{m} > 0,
\]
und folglich erhalten wir mit R\"ucksicht auf (12), (13), wenn wir
$\Theta^{-1}$ f\"ur $\Theta$ setzen, wodurch $\chi(t)^{-1}$ in $\chi(t)$
\"ubergeht:
\weq{15}{X_1 = -\frac{2}{m} (+1, \Theta^{-1}, r) \sum^{t} \chi_1(t)
\log\left(\operatorname{tg} \frac{t\pi}{m}\right),}
\weq{16}{X_2 = \frac{i\pi}{m} (-1, \Theta^{-1}, r) \sum^{t} \chi_2(t),
\qquad 1 \leqq t \leqq \nu - 1.}

\begin{center}
\S.\,198.

\gesperrt{Der Classenzahlfactor $A$.}
\end{center}

Um nun zun\"achst $A$ zu berechnen, haben wir den Ausdruck (16) f\"ur
$X_2$ in die erste Formel (8) des vorigen Paragraphen einzusetzen.
Diese Formel k\"onnen wir auch so darstellen:
\weq{1}{\pi^{\nu'} A = 2\, m^{1/4\,\nu} 2^{-\nu'} \prod X_2.}
